\clearpage
\section{Apendice: mapa pagina por pagina de los slides de integer}
Este apendice lista cada pagina/slide de la carpeta fuente y su contenido principal. Muchas paginas son animaciones repetidas de Beamer; se conservan para trazabilidad.
\subsection*{Estructuras\_de\_Datos\_Avanzadas\_\_\_Modelo\_de\_enteros\_II.pdf (52 paginas)}
\begin{description}[leftmargin=!,labelwidth=4.2em,style=nextline]
\item[Pag. 1] CS3014 Estructuras de Datos Avanzadas CS3014 Estructuras de Datos Avanzadas Laboratorio - Semana 9 - Sesion 1 Victor Racso Galvan Oyola vgalvan@utec.edu.pe 20 de mayo de 2026 Victor Racso Galvan Oyola Universidad de Ingenieria y Tecnologia 1 / 14
\item[Pag. 2] CS3014 Estructuras de Datos Avanzadas ?Que aprenderemos hoy? -Modelos computacional RAM AC 0 -Fusion Trees Victor Racso Galvan Oyola Universidad de Ingenieria y Tecnologia 2 / 14
\item[Pag. 3] CS3014 Estructuras de Datos Avanzadas ?Que aprenderemos hoy? -Modelos computacional RAM AC 0 -Fusion Trees Victor Racso Galvan Oyola Universidad de Ingenieria y Tecnologia 2 / 14
\item[Pag. 4] CS3014 Estructuras de Datos Avanzadas ?Que aprenderemos hoy? -Modelos computacional RAM AC 0 -Fusion Trees Victor Racso Galvan Oyola Universidad de Ingenieria y Tecnologia 2 / 14
\item[Pag. 5] Modelos computacional RAM AC 0
\item[Pag. 6] CS3014 Estructuras de Datos Avanzadas Modelos computacional RAM AC 0 Modelo RAM AC0 -Permite solo operaciones que se puedan realizar en un circuito de profundidadO(1)pero sin limite de fan-in y fan-out (compuertas que entran y salen, respectivamente). -Por su definicion, no permite multiplicacion. Otras operaciones permitidas Operaciones bitwise pueden ser realizadas enO(1). Victor Racso Galvan Oyola Universidad de Ingenieria y Tecnologia 4 / 14
\item[Pag. 7] CS3014 Estructuras de Datos Avanzadas Modelos computacional RAM AC 0 Modelo RAM AC0 -Permite solo operaciones que se puedan realizar en un circuito de profundidadO(1)pero sin limite de fan-in y fan-out (compuertas que entran y salen, respectivamente). -Por su definicion, no permite multiplicacion. Otras operaciones permitidas Operaciones bitwise pueden ser realizadas enO(1). Victor Racso Galvan Oyola Universidad de Ingenieria y Tecnologia 4 / 14
\item[Pag. 8] CS3014 Estructuras de Datos Avanzadas Modelos computacional RAM AC 0 Modelo RAM AC0 -Permite solo operaciones que se puedan realizar en un circuito de profundidadO(1)pero sin limite de fan-in y fan-out (compuertas que entran y salen, respectivamente). -Por su definicion, no permite multiplicacion. Otras operaciones permitidas Operaciones bitwise pueden ser realizadas enO(1). Victor Racso Galvan Oyola Universidad de Ingenieria y Tecnologia 4 / 14
\item[Pag. 9] CS3014 Estructuras de Datos Avanzadas Modelos computacional RAM AC 0 Modelo RAM AC0 -Permite solo operaciones que se puedan realizar en un circuito de profundidadO(1)pero sin limite de fan-in y fan-out (compuertas que entran y salen, respectivamente). -Por su definicion, no permite multiplicacion. Otras operaciones permitidas Operaciones bitwise pueden ser realizadas enO(1). Victor Racso Galvan Oyola Universidad de Ingenieria y Tecnologia 4 / 14
\item[Pag. 10] CS3014 Estructuras de Datos Avanzadas Modelos computacional RAM AC 0 Modelo RAM AC0 -Permite solo operaciones que se puedan realizar en un circuito de profundidadO(1)pero sin limite de fan-in y fan-out (compuertas que entran y salen, respectivamente). -Por su definicion, no permite multiplicacion. Otras operaciones permitidas Operaciones bitwise pueden ser realizadas enO(1). Victor Racso Galvan Oyola Universidad de Ingenieria y Tecnologia 4 / 14
\item[Pag. 11] Fusion T rees
\item[Pag. 12] CS3014 Estructuras de Datos Avanzadas Fusion Trees Definicion Es una estructura de datos que nos permite consultar el sucesor/predecesor enO(log w n). Considerando este resultado y el de vEB y y-fast trees, tendriamos que la consulta de sucesor predecesor toma m ' in\{O(logw),O(logw n)\}=O( p logn) Versiones Veremos la version estatica. Las dinamicas disponibles cumplen: -O(log w n+ log logn)usando exponential trees. -O(log w n)esperado usando hashing. Victor Racso Galvan Oyola Universidad de Ingenieria y Tecnologia 6 / 14
\item[Pag. 13] CS3014 Estructuras de Datos Avanzadas Fusion Trees Definicion Es una estructura de datos que nos permite consultar el sucesor/predecesor enO(log w n). Considerando este resultado y el de vEB y y-fast trees, tendriamos que la consulta de sucesor predecesor toma m ' in\{O(logw),O(logw n)\}=O( p logn) Versiones Veremos la version estatica. Las dinamicas disponibles cumplen: -O(log w n+ log logn)usando exponential trees. -O(log w n)esperado usando hashing. Victor Racso Galvan Oyola Universidad de Ingenieria y Tecnologia 6 / 14
\item[Pag. 14] CS3014 Estructuras de Datos Avanzadas Fusion Trees Definicion Es una estructura de datos que nos permite consultar el sucesor/predecesor enO(log w n). Considerando este resultado y el de vEB y y-fast trees, tendriamos que la consulta de sucesor predecesor toma m ' in\{O(logw),O(logw n)\}=O( p logn) Versiones Veremos la version estatica. Las dinamicas disponibles cumplen: -O(log w n+ log logn)usando exponential trees. -O(log w n)esperado usando hashing. Victor Racso Galvan Oyola Universidad de Ingenieria y Tecnologia 6 / 14
\item[Pag. 15] CS3014 Estructuras de Datos Avanzadas Fusion Trees Definicion Es una estructura de datos que nos permite consultar el sucesor/predecesor enO(log w n). Considerando este resultado y el de vEB y y-fast trees, tendriamos que la consulta de sucesor predecesor toma m ' in\{O(logw),O(logw n)\}=O( p logn) Versiones Veremos la version estatica. Las dinamicas disponibles cumplen: -O(log w n+ log logn)usando exponential trees. -O(log w n)esperado usando hashing. Victor Racso Galvan Oyola Universidad de Ingenieria y Tecnologia 6 / 14
\item[Pag. 16] CS3014 Estructuras de Datos Avanzadas Fusion Trees Definicion Es una estructura de datos que nos permite consultar el sucesor/predecesor enO(log w n). Considerando este resultado y el de vEB y y-fast trees, tendriamos que la consulta de sucesor predecesor toma m ' in\{O(logw),O(logw n)\}=O( p logn) Versiones Veremos la version estatica. Las dinamicas disponibles cumplen: -O(log w n+ log logn)usando exponential trees. -O(log w n)esperado usando hashing. Victor Racso Galvan Oyola Universidad de Ingenieria y Tecnologia 6 / 14
\item[Pag. 17] CS3014 Estructuras de Datos Avanzadas Fusion Trees Definicion Es una estructura de datos que nos permite consultar el sucesor/predecesor enO(log w n). Considerando este resultado y el de vEB y y-fast trees, tendriamos que la consulta de sucesor predecesor toma m ' in\{O(logw),O(logw n)\}=O( p logn) Versiones Veremos la version estatica. Las dinamicas disponibles cumplen: -O(log w n+ log logn)usando exponential trees. -O(log w n)esperado usando hashing. Victor Racso Galvan Oyola Universidad de Ingenieria y Tecnologia 6 / 14
\item[Pag. 18] CS3014 Estructuras de Datos Avanzadas Fusion Trees Definicion Es una estructura de datos que nos permite consultar el sucesor/predecesor enO(log w n). Considerando este resultado y el de vEB y y-fast trees, tendriamos que la consulta de sucesor predecesor toma m ' in\{O(logw),O(logw n)\}=O( p logn) Versiones Veremos la version estatica. Las dinamicas disponibles cumplen: -O(log w n+ log logn)usando exponential trees. -O(log w n)esperado usando hashing. Victor Racso Galvan Oyola Universidad de Ingenieria y Tecnologia 6 / 14
\item[Pag. 19] CS3014 Estructuras de Datos Avanzadas Estructura Planteamiento inicial Consideraremos un B-Tree de ordenw 1 5, asi que cada nodo tiene a lo mucho w 1 5 hijos. Cada nodo manejaw 1+ 1 5 bits, asi que no podemos buscar el hijo correspon- diente enO(1). Solucion - Nodos de fusion tree -Almacenamosk=O(w 1 5 )valoresx 0 <x 1 < . . . <x k-1. -Debemos poder responder al sucesor o predecesor de un valor enO(1). -Requeriremosk O(1) de preprocesamiento. Victor Racso Galvan Oyola Universidad de Ingenieria y Tecnologia 7 / 14
\item[Pag. 20] CS3014 Estructuras de Datos Avanzadas Estructura Planteamiento inicial Consideraremos un B-Tree de ordenw 1 5, asi que cada nodo tiene a lo mucho w 1 5 hijos. Cada nodo manejaw 1+ 1 5 bits, asi que no podemos buscar el hijo correspon- diente enO(1). Solucion - Nodos de fusion tree -Almacenamosk=O(w 1 5 )valoresx 0 <x 1 < . . . <x k-1. -Debemos poder responder al sucesor o predecesor de un valor enO(1). -Requeriremosk O(1) de preprocesamiento. Victor Racso Galvan Oyola Universidad de Ingenieria y Tecnologia 7 / 14
\item[Pag. 21] CS3014 Estructuras de Datos Avanzadas Estructura Planteamiento inicial Consideraremos un B-Tree de ordenw 1 5, asi que cada nodo tiene a lo mucho w 1 5 hijos. Cada nodo manejaw 1+ 1 5 bits, asi que no podemos buscar el hijo correspon- diente enO(1). Solucion - Nodos de fusion tree -Almacenamosk=O(w 1 5 )valoresx 0 <x 1 < . . . <x k-1. -Debemos poder responder al sucesor o predecesor de un valor enO(1). -Requeriremosk O(1) de preprocesamiento. Victor Racso Galvan Oyola Universidad de Ingenieria y Tecnologia 7 / 14
\item[Pag. 22] CS3014 Estructuras de Datos Avanzadas Estructura Planteamiento inicial Consideraremos un B-Tree de ordenw 1 5, asi que cada nodo tiene a lo mucho w 1 5 hijos. Cada nodo manejaw 1+ 1 5 bits, asi que no podemos buscar el hijo correspon- diente enO(1). Solucion - Nodos de fusion tree -Almacenamosk=O(w 1 5 )valoresx 0 <x 1 < . . . <x k-1. -Debemos poder responder al sucesor o predecesor de un valor enO(1). -Requeriremosk O(1) de preprocesamiento. Victor Racso Galvan Oyola Universidad de Ingenieria y Tecnologia 7 / 14
\item[Pag. 23] CS3014 Estructuras de Datos Avanzadas Estructura Planteamiento inicial Consideraremos un B-Tree de ordenw 1 5, asi que cada nodo tiene a lo mucho w 1 5 hijos. Cada nodo manejaw 1+ 1 5 bits, asi que no podemos buscar el hijo correspon- diente enO(1). Solucion - Nodos de fusion tree -Almacenamosk=O(w 1 5 )valoresx 0 <x 1 < . . . <x k-1. -Debemos poder responder al sucesor o predecesor de un valor enO(1). -Requeriremosk O(1) de preprocesamiento. Victor Racso Galvan Oyola Universidad de Ingenieria y Tecnologia 7 / 14
\item[Pag. 24] CS3014 Estructuras de Datos Avanzadas Estructura Planteamiento inicial Consideraremos un B-Tree de ordenw 1 5, asi que cada nodo tiene a lo mucho w 1 5 hijos. Cada nodo manejaw 1+ 1 5 bits, asi que no podemos buscar el hijo correspon- diente enO(1). Solucion - Nodos de fusion tree -Almacenamosk=O(w 1 5 )valoresx 0 <x 1 < . . . <x k-1. -Debemos poder responder al sucesor o predecesor de un valor enO(1). -Requeriremosk O(1) de preprocesamiento. Victor Racso Galvan Oyola Universidad de Ingenieria y Tecnologia 7 / 14
\item[Pag. 25] CS3014 Estructuras de Datos Avanzadas Estructura Planteamiento inicial Consideraremos un B-Tree de ordenw 1 5, asi que cada nodo tiene a lo mucho w 1 5 hijos. Cada nodo manejaw 1+ 1 5 bits, asi que no podemos buscar el hijo correspon- diente enO(1). Solucion - Nodos de fusion tree -Almacenamosk=O(w 1 5 )valoresx 0 <x 1 < . . . <x k-1. -Debemos poder responder al sucesor o predecesor de un valor enO(1). -Requeriremosk O(1) de preprocesamiento. Victor Racso Galvan Oyola Universidad de Ingenieria y Tecnologia 7 / 14
\item[Pag. 26] CS3014 Estructuras de Datos Avanzadas Sketching ?Como distinguimosk=O(w 1 5 )valores? Podemos considerar un trie de las llaves con un alfabeto binario. Es posible reducir la cantidad total de nodos calculando el auxiliary/virtual tree del trie correspondiente y mapeamos los niveles involucrados. Consecuencias Hayk-1 nodos internos en el trie, asi que nos interesan a lo muchok-1 niveles del mismo. Si almacenamos la informacion de esosk-1 niveles, entonces manejaremos O(w 1 5 )bits por cada llave. En total, manejaremosO(w 2 5 )bits, lo cual si es posible d...
\item[Pag. 27] CS3014 Estructuras de Datos Avanzadas Sketching ?Como distinguimosk=O(w 1 5 )valores? Podemos considerar un trie de las llaves con un alfabeto binario. Es posible reducir la cantidad total de nodos calculando el auxiliary/virtual tree del trie correspondiente y mapeamos los niveles involucrados. Consecuencias Hayk-1 nodos internos en el trie, asi que nos interesan a lo muchok-1 niveles del mismo. Si almacenamos la informacion de esosk-1 niveles, entonces manejaremos O(w 1 5 )bits por cada llave. En total, manejaremosO(w 2 5 )bits, lo cual si es posible d...
\item[Pag. 28] CS3014 Estructuras de Datos Avanzadas Sketching ?Como distinguimosk=O(w 1 5 )valores? Podemos considerar un trie de las llaves con un alfabeto binario. Es posible reducir la cantidad total de nodos calculando el auxiliary/virtual tree del trie correspondiente y mapeamos los niveles involucrados. Consecuencias Hayk-1 nodos internos en el trie, asi que nos interesan a lo muchok-1 niveles del mismo. Si almacenamos la informacion de esosk-1 niveles, entonces manejaremos O(w 1 5 )bits por cada llave. En total, manejaremosO(w 2 5 )bits, lo cual si es posible d...
\item[Pag. 29] CS3014 Estructuras de Datos Avanzadas Sketching ?Como distinguimosk=O(w 1 5 )valores? Podemos considerar un trie de las llaves con un alfabeto binario. Es posible reducir la cantidad total de nodos calculando el auxiliary/virtual tree del trie correspondiente y mapeamos los niveles involucrados. Consecuencias Hayk-1 nodos internos en el trie, asi que nos interesan a lo muchok-1 niveles del mismo. Si almacenamos la informacion de esosk-1 niveles, entonces manejaremos O(w 1 5 )bits por cada llave. En total, manejaremosO(w 2 5 )bits, lo cual si es posible d...
\item[Pag. 30] CS3014 Estructuras de Datos Avanzadas Sketching ?Como distinguimosk=O(w 1 5 )valores? Podemos considerar un trie de las llaves con un alfabeto binario. Es posible reducir la cantidad total de nodos calculando el auxiliary/virtual tree del trie correspondiente y mapeamos los niveles involucrados. Consecuencias Hayk-1 nodos internos en el trie, asi que nos interesan a lo muchok-1 niveles del mismo. Si almacenamos la informacion de esosk-1 niveles, entonces manejaremos O(w 1 5 )bits por cada llave. En total, manejaremosO(w 2 5 )bits, lo cual si es posible d...
\item[Pag. 31] CS3014 Estructuras de Datos Avanzadas Sketching ?Como distinguimosk=O(w 1 5 )valores? Podemos considerar un trie de las llaves con un alfabeto binario. Es posible reducir la cantidad total de nodos calculando el auxiliary/virtual tree del trie correspondiente y mapeamos los niveles involucrados. Consecuencias Hayk-1 nodos internos en el trie, asi que nos interesan a lo muchok-1 niveles del mismo. Si almacenamos la informacion de esosk-1 niveles, entonces manejaremos O(w 1 5 )bits por cada llave. En total, manejaremosO(w 2 5 )bits, lo cual si es posible d...
\item[Pag. 32] CS3014 Estructuras de Datos Avanzadas Sketching ?Como distinguimosk=O(w 1 5 )valores? Podemos considerar un trie de las llaves con un alfabeto binario. Es posible reducir la cantidad total de nodos calculando el auxiliary/virtual tree del trie correspondiente y mapeamos los niveles involucrados. Consecuencias Hayk-1 nodos internos en el trie, asi que nos interesan a lo muchok-1 niveles del mismo. Si almacenamos la informacion de esosk-1 niveles, entonces manejaremos O(w 1 5 )bits por cada llave. En total, manejaremosO(w 2 5 )bits, lo cual si es posible d...
\item[Pag. 33] CS3014 Estructuras de Datos Avanzadas Trie binario empty 1 1 0 2 0 3 1 1 0 5 1 0 1 1=001,2=010,3=011,5=101 Victor Racso Galvan Oyola Universidad de Ingenieria y Tecnologia 9 / 14
\item[Pag. 34] CS3014 Estructuras de Datos Avanzadas Trie binario + Virtual Tree empty 1 1 0 2 0 3 1 1 0 5 1 0 1 1=001,2=010,3=011,5=101 Victor Racso Galvan Oyola Universidad de Ingenieria y Tecnologia 10 / 14
\item[Pag. 35] CS3014 Estructuras de Datos Avanzadas Sketching perfecto Definicion Consideraremos una secuencia derenterosb 0 >b 1 > . . . >b r-1 y cada llave la representaremos como una cadena derbits en la cual eli-esimo bit correspondera al bitbi de la llave. Denotaremos a dicha cadena binaria comoSketch(x). Propiedad importante Sketch(x0)<Sketch(x 1)< . . . <Sketch(x k-1) Esto se cumplesolo para las llaves. Victor Racso Galvan Oyola Universidad de Ingenieria y Tecnologia 11 / 14
\item[Pag. 36] CS3014 Estructuras de Datos Avanzadas Sketching perfecto Definicion Consideraremos una secuencia derenterosb 0 >b 1 > . . . >b r-1 y cada llave la representaremos como una cadena derbits en la cual eli-esimo bit correspondera al bitbi de la llave. Denotaremos a dicha cadena binaria comoSketch(x). Propiedad importante Sketch(x0)<Sketch(x 1)< . . . <Sketch(x k-1) Esto se cumplesolo para las llaves. Victor Racso Galvan Oyola Universidad de Ingenieria y Tecnologia 11 / 14
\item[Pag. 37] CS3014 Estructuras de Datos Avanzadas Sketching perfecto Definicion Consideraremos una secuencia derenterosb 0 >b 1 > . . . >b r-1 y cada llave la representaremos como una cadena derbits en la cual eli-esimo bit correspondera al bitbi de la llave. Denotaremos a dicha cadena binaria comoSketch(x). Propiedad importante Sketch(x0)<Sketch(x 1)< . . . <Sketch(x k-1) Esto se cumplesolo para las llaves. Victor Racso Galvan Oyola Universidad de Ingenieria y Tecnologia 11 / 14
\item[Pag. 38] CS3014 Estructuras de Datos Avanzadas Sketching perfecto Definicion Consideraremos una secuencia derenterosb 0 >b 1 > . . . >b r-1 y cada llave la representaremos como una cadena derbits en la cual eli-esimo bit correspondera al bitbi de la llave. Denotaremos a dicha cadena binaria comoSketch(x). Propiedad importante Sketch(x0)<Sketch(x 1)< . . . <Sketch(x k-1) Esto se cumplesolo para las llaves. Victor Racso Galvan Oyola Universidad de Ingenieria y Tecnologia 11 / 14
\item[Pag. 39] CS3014 Estructuras de Datos Avanzadas Sketching perfecto Definicion Consideraremos una secuencia derenterosb 0 >b 1 > . . . >b r-1 y cada llave la representaremos como una cadena derbits en la cual eli-esimo bit correspondera al bitbi de la llave. Denotaremos a dicha cadena binaria comoSketch(x). Propiedad importante Sketch(x0)<Sketch(x 1)< . . . <Sketch(x k-1) Esto se cumplesolo para las llaves. Victor Racso Galvan Oyola Universidad de Ingenieria y Tecnologia 11 / 14
\item[Pag. 40] CS3014 Estructuras de Datos Avanzadas Sketching perfecto Definicion Consideraremos una secuencia derenterosb 0 >b 1 > . . . >b r-1 y cada llave la representaremos como una cadena derbits en la cual eli-esimo bit correspondera al bitbi de la llave. Denotaremos a dicha cadena binaria comoSketch(x). Propiedad importante Sketch(x0)<Sketch(x 1)< . . . <Sketch(x k-1) Esto se cumplesolo para las llaves. Victor Racso Galvan Oyola Universidad de Ingenieria y Tecnologia 11 / 14
\item[Pag. 41] CS3014 Estructuras de Datos Avanzadas Consultas ?Como buscamos el sucesor o predecesor de un valorq? Podemos pensar en calcularSketch(q)y luego buscarlo en la estructura. Sorpresa:Sketchpreserva el orden entre las llaves, asi queSketch(q) podria caer en cualquier lugar, no necesariamente en su sucesor o predecesor. ?Que hacemos? Supongamos queSketch(x i)<=Sketch(q)<Sketch(x i+1). Tomamos ellongest common prefixentreqyx i ox i+1. El nodo al que lleguemos es en donde la busquedafalla, asi que podremos decidir que hacer a continuacion dependiendo del escena...
\item[Pag. 42] CS3014 Estructuras de Datos Avanzadas Consultas ?Como buscamos el sucesor o predecesor de un valorq? Podemos pensar en calcularSketch(q)y luego buscarlo en la estructura. Sorpresa:Sketchpreserva el orden entre las llaves, asi queSketch(q) podria caer en cualquier lugar, no necesariamente en su sucesor o predecesor. ?Que hacemos? Supongamos queSketch(x i)<=Sketch(q)<Sketch(x i+1). Tomamos ellongest common prefixentreqyx i ox i+1. El nodo al que lleguemos es en donde la busquedafalla, asi que podremos decidir que hacer a continuacion dependiendo del escena...
\item[Pag. 43] CS3014 Estructuras de Datos Avanzadas Consultas ?Como buscamos el sucesor o predecesor de un valorq? Podemos pensar en calcularSketch(q)y luego buscarlo en la estructura. Sorpresa:Sketchpreserva el orden entre las llaves, asi queSketch(q) podria caer en cualquier lugar, no necesariamente en su sucesor o predecesor. ?Que hacemos? Supongamos queSketch(x i)<=Sketch(q)<Sketch(x i+1). Tomamos ellongest common prefixentreqyx i ox i+1. El nodo al que lleguemos es en donde la busquedafalla, asi que podremos decidir que hacer a continuacion dependiendo del escena...
\item[Pag. 44] CS3014 Estructuras de Datos Avanzadas Consultas ?Como buscamos el sucesor o predecesor de un valorq? Podemos pensar en calcularSketch(q)y luego buscarlo en la estructura. Sorpresa:Sketchpreserva el orden entre las llaves, asi queSketch(q) podria caer en cualquier lugar, no necesariamente en su sucesor o predecesor. ?Que hacemos? Supongamos queSketch(x i)<=Sketch(q)<Sketch(x i+1). Tomamos ellongest common prefixentreqyx i ox i+1. El nodo al que lleguemos es en donde la busquedafalla, asi que podremos decidir que hacer a continuacion dependiendo del escena...
\item[Pag. 45] CS3014 Estructuras de Datos Avanzadas Consultas ?Como buscamos el sucesor o predecesor de un valorq? Podemos pensar en calcularSketch(q)y luego buscarlo en la estructura. Sorpresa:Sketchpreserva el orden entre las llaves, asi queSketch(q) podria caer en cualquier lugar, no necesariamente en su sucesor o predecesor. ?Que hacemos? Supongamos queSketch(x i)<=Sketch(q)<Sketch(x i+1). Tomamos ellongest common prefixentreqyx i ox i+1. El nodo al que lleguemos es en donde la busquedafalla, asi que podremos decidir que hacer a continuacion dependiendo del escena...
\item[Pag. 46] CS3014 Estructuras de Datos Avanzadas Consultas ?Como buscamos el sucesor o predecesor de un valorq? Podemos pensar en calcularSketch(q)y luego buscarlo en la estructura. Sorpresa:Sketchpreserva el orden entre las llaves, asi queSketch(q) podria caer en cualquier lugar, no necesariamente en su sucesor o predecesor. ?Que hacemos? Supongamos queSketch(x i)<=Sketch(q)<Sketch(x i+1). Tomamos ellongest common prefixentreqyx i ox i+1. El nodo al que lleguemos es en donde la busquedafalla, asi que podremos decidir que hacer a continuacion dependiendo del escena...
\item[Pag. 47] CS3014 Estructuras de Datos Avanzadas Consultas ?Como buscamos el sucesor o predecesor de un valorq? Podemos pensar en calcularSketch(q)y luego buscarlo en la estructura. Sorpresa:Sketchpreserva el orden entre las llaves, asi queSketch(q) podria caer en cualquier lugar, no necesariamente en su sucesor o predecesor. ?Que hacemos? Supongamos queSketch(x i)<=Sketch(q)<Sketch(x i+1). Tomamos ellongest common prefixentreqyx i ox i+1. El nodo al que lleguemos es en donde la busquedafalla, asi que podremos decidir que hacer a continuacion dependiendo del escena...
\item[Pag. 48] CS3014 Estructuras de Datos Avanzadas Consultas ?Como buscamos el sucesor o predecesor de un valorq? Podemos pensar en calcularSketch(q)y luego buscarlo en la estructura. Sorpresa:Sketchpreserva el orden entre las llaves, asi queSketch(q) podria caer en cualquier lugar, no necesariamente en su sucesor o predecesor. ?Que hacemos? Supongamos queSketch(x i)<=Sketch(q)<Sketch(x i+1). Tomamos ellongest common prefixentreqyx i ox i+1. El nodo al que lleguemos es en donde la busquedafalla, asi que podremos decidir que hacer a continuacion dependiendo del escena...
\item[Pag. 49] CS3014 Estructuras de Datos Avanzadas Resumen de la sesion -Aprendimos un poco del modelo RAM AC0 . -Aprendimos sobre Fusion Trees. Victor Racso Galvan Oyola Universidad de Ingenieria y Tecnologia 13 / 14
\item[Pag. 50] CS3014 Estructuras de Datos Avanzadas Resumen de la sesion -Aprendimos un poco del modelo RAM AC0 . -Aprendimos sobre Fusion Trees. Victor Racso Galvan Oyola Universidad de Ingenieria y Tecnologia 13 / 14
\item[Pag. 51] CS3014 Estructuras de Datos Avanzadas Resumen de la sesion -Aprendimos un poco del modelo RAM AC0 . -Aprendimos sobre Fusion Trees. Victor Racso Galvan Oyola Universidad de Ingenieria y Tecnologia 13 / 14
\item[Pag. 52] CS3014 Estructuras de Datos Avanzadas Gracias Victor Racso Galvan Oyola Universidad de Ingenieria y Tecnologia 14 / 14
\end{description}
\subsection*{Estructuras\_de\_Datos\_Avanzadas\_\_\_Ordenamiento\_en\_tiempo\_lineal.pdf (59 paginas)}
\begin{description}[leftmargin=!,labelwidth=4.2em,style=nextline]
\item[Pag. 1] CS3014 Estructuras de Datos Avanzadas CS3014 Estructuras de Datos Avanzadas Laboratorio - Semana 9 - Sesion 2 Victor Racso Galvan Oyola vgalvan@utec.edu.pe 22 de mayo de 2026 Victor Racso Galvan Oyola Universidad de Ingenieria y Tecnologia 1 / 15
\item[Pag. 2] CS3014 Estructuras de Datos Avanzadas ?Que aprenderemos hoy? -Equivalencia entre ordenamiento y colas de prioridad -Algunos resultados en ordenamiento de enteros -Signature sort Victor Racso Galvan Oyola Universidad de Ingenieria y Tecnologia 2 / 15
\item[Pag. 3] CS3014 Estructuras de Datos Avanzadas ?Que aprenderemos hoy? -Equivalencia entre ordenamiento y colas de prioridad -Algunos resultados en ordenamiento de enteros -Signature sort Victor Racso Galvan Oyola Universidad de Ingenieria y Tecnologia 2 / 15
\item[Pag. 4] CS3014 Estructuras de Datos Avanzadas ?Que aprenderemos hoy? -Equivalencia entre ordenamiento y colas de prioridad -Algunos resultados en ordenamiento de enteros -Signature sort Victor Racso Galvan Oyola Universidad de Ingenieria y Tecnologia 2 / 15
\item[Pag. 5] CS3014 Estructuras de Datos Avanzadas ?Que aprenderemos hoy? -Equivalencia entre ordenamiento y colas de prioridad -Algunos resultados en ordenamiento de enteros -Signature sort Victor Racso Galvan Oyola Universidad de Ingenieria y Tecnologia 2 / 15
\item[Pag. 6] Ordenamiento y cola de prioridades
\item[Pag. 7] CS3014 Estructuras de Datos Avanzadas Equivalencia Punto inicial Es sencillo notar que si tenemos una cola de prioridad que soporta las opera- ciones clasicas enO(f(n,w)), entonces existe un algoritmo de ordenamiento que tomaO(nf(n,w)). Punto final Si existe un algoritmo de ordenamiento que tomaO(T(n,w)), existe una cola de prioridad que soporta sus operaciones clasicas enO T(n,w) n . Ejemplo: Quickheap. Victor Racso Galvan Oyola Universidad de Ingenieria y Tecnologia 4 / 15
\item[Pag. 8] CS3014 Estructuras de Datos Avanzadas Equivalencia Punto inicial Es sencillo notar que si tenemos una cola de prioridad que soporta las opera- ciones clasicas enO(f(n,w)), entonces existe un algoritmo de ordenamiento que tomaO(nf(n,w)). Punto final Si existe un algoritmo de ordenamiento que tomaO(T(n,w)), existe una cola de prioridad que soporta sus operaciones clasicas enO T(n,w) n . Ejemplo: Quickheap. Victor Racso Galvan Oyola Universidad de Ingenieria y Tecnologia 4 / 15
\item[Pag. 9] CS3014 Estructuras de Datos Avanzadas Equivalencia Punto inicial Es sencillo notar que si tenemos una cola de prioridad que soporta las opera- ciones clasicas enO(f(n,w)), entonces existe un algoritmo de ordenamiento que tomaO(nf(n,w)). Punto final Si existe un algoritmo de ordenamiento que tomaO(T(n,w)), existe una cola de prioridad que soporta sus operaciones clasicas enO T(n,w) n . Ejemplo: Quickheap. Victor Racso Galvan Oyola Universidad de Ingenieria y Tecnologia 4 / 15
\item[Pag. 10] CS3014 Estructuras de Datos Avanzadas Equivalencia Punto inicial Es sencillo notar que si tenemos una cola de prioridad que soporta las opera- ciones clasicas enO(f(n,w)), entonces existe un algoritmo de ordenamiento que tomaO(nf(n,w)). Punto final Si existe un algoritmo de ordenamiento que tomaO(T(n,w)), existe una cola de prioridad que soporta sus operaciones clasicas enO T(n,w) n . Ejemplo: Quickheap. Victor Racso Galvan Oyola Universidad de Ingenieria y Tecnologia 4 / 15
\item[Pag. 11] CS3014 Estructuras de Datos Avanzadas Equivalencia Punto inicial Es sencillo notar que si tenemos una cola de prioridad que soporta las opera- ciones clasicas enO(f(n,w)), entonces existe un algoritmo de ordenamiento que tomaO(nf(n,w)). Punto final Si existe un algoritmo de ordenamiento que tomaO(T(n,w)), existe una cola de prioridad que soporta sus operaciones clasicas enO T(n,w) n . Ejemplo: Quickheap. Victor Racso Galvan Oyola Universidad de Ingenieria y Tecnologia 4 / 15
\item[Pag. 12] Resultados en ordenamientos de enteros
\item[Pag. 13] CS3014 Estructuras de Datos Avanzadas Resultados en ordenamiento de enteros Para ordenarnenteros dewbits, la siguiente tabla resume algunos resultados historicos. Metodo Complejidad Condicion Comparaciones O(nlogn) Sin condiciones Counting sort O(n+2 w ) O(n)siw= logn Radix sort O n w logn O(n)siw=O(logn) van Emde Boas O(nlogw) Basado en van Emde Boas O(nlog w logn ) O(nlog logn)siw= log O(1) n Signature sort O(n) w= Omega(log 2+epsilon n), forall epsilon >0 Han (2004) O(nlog logn) modelo RAM AC0 Han y Thorup (2002) O(n q log w logn ) Aleatorizado Cuadro...
\item[Pag. 14] CS3014 Estructuras de Datos Avanzadas Resultados en ordenamiento de enteros Para ordenarnenteros dewbits, la siguiente tabla resume algunos resultados historicos. Metodo Complejidad Condicion Comparaciones O(nlogn) Sin condiciones Counting sort O(n+2 w ) O(n)siw= logn Radix sort O n w logn O(n)siw=O(logn) van Emde Boas O(nlogw) Basado en van Emde Boas O(nlog w logn ) O(nlog logn)siw= log O(1) n Signature sort O(n) w= Omega(log 2+epsilon n), forall epsilon >0 Han (2004) O(nlog logn) modelo RAM AC0 Han y Thorup (2002) O(n q log w logn ) Aleatorizado Cuadro...
\item[Pag. 15] CS3014 Estructuras de Datos Avanzadas Resultados en ordenamiento de enteros Para ordenarnenteros dewbits, la siguiente tabla resume algunos resultados historicos. Metodo Complejidad Condicion Comparaciones O(nlogn) Sin condiciones Counting sort O(n+2 w ) O(n)siw= logn Radix sort O n w logn O(n)siw=O(logn) van Emde Boas O(nlogw) Basado en van Emde Boas O(nlog w logn ) O(nlog logn)siw= log O(1) n Signature sort O(n) w= Omega(log 2+epsilon n), forall epsilon >0 Han (2004) O(nlog logn) modelo RAM AC0 Han y Thorup (2002) O(n q log w logn ) Aleatorizado Cuadro...
\item[Pag. 16] Signature sort
\item[Pag. 17] CS3014 Estructuras de Datos Avanzadas Signature sort Condicion base Asumiremos quew>=log 2+epsilon nlog lognpor ser conveniente. Paso 1 Particionamos cada entero enlogepsilon nbloques de tamanos iguales que seran digitos. Notemos que estos bloques tendran tamano de al menoslog2 n. Esto es equivalente a tomar una baseB\textasciitilde{} w logepsilon n y representar cada entero en baseB. Victor Racso Galvan Oyola Universidad de Ingenieria y Tecnologia 8 / 15
\item[Pag. 18] CS3014 Estructuras de Datos Avanzadas Signature sort Condicion base Asumiremos quew>=log 2+epsilon nlog lognpor ser conveniente. Paso 1 Particionamos cada entero enlogepsilon nbloques de tamanos iguales que seran digitos. Notemos que estos bloques tendran tamano de al menoslog2 n. Esto es equivalente a tomar una baseB\textasciitilde{} w logepsilon n y representar cada entero en baseB. Victor Racso Galvan Oyola Universidad de Ingenieria y Tecnologia 8 / 15
\item[Pag. 19] CS3014 Estructuras de Datos Avanzadas Signature sort Condicion base Asumiremos quew>=log 2+epsilon nlog lognpor ser conveniente. Paso 1 Particionamos cada entero enlogepsilon nbloques de tamanos iguales que seran digitos. Notemos que estos bloques tendran tamano de al menoslog2 n. Esto es equivalente a tomar una baseB\textasciitilde{} w logepsilon n y representar cada entero en baseB. Victor Racso Galvan Oyola Universidad de Ingenieria y Tecnologia 8 / 15
\item[Pag. 20] CS3014 Estructuras de Datos Avanzadas Signature sort Condicion base Asumiremos quew>=log 2+epsilon nlog lognpor ser conveniente. Paso 1 Particionamos cada entero enlogepsilon nbloques de tamanos iguales que seran digitos. Notemos que estos bloques tendran tamano de al menoslog2 n. Esto es equivalente a tomar una baseB\textasciitilde{} w logepsilon n y representar cada entero en baseB. Victor Racso Galvan Oyola Universidad de Ingenieria y Tecnologia 8 / 15
\item[Pag. 21] CS3014 Estructuras de Datos Avanzadas Signature sort Condicion base Asumiremos quew>=log 2+epsilon nlog lognpor ser conveniente. Paso 1 Particionamos cada entero enlogepsilon nbloques de tamanos iguales que seran digitos. Notemos que estos bloques tendran tamano de al menoslog2 n. Esto es equivalente a tomar una baseB\textasciitilde{} w logepsilon n y representar cada entero en baseB. Victor Racso Galvan Oyola Universidad de Ingenieria y Tecnologia 8 / 15
\item[Pag. 22] CS3014 Estructuras de Datos Avanzadas Signature sort Condicion base Asumiremos quew>=log 2+epsilon nlog lognpor ser conveniente. Paso 1 Particionamos cada entero enlogepsilon nbloques de tamanos iguales que seran digitos. Notemos que estos bloques tendran tamano de al menoslog2 n. Esto es equivalente a tomar una baseB\textasciitilde{} w logepsilon n y representar cada entero en baseB. Victor Racso Galvan Oyola Universidad de Ingenieria y Tecnologia 8 / 15
\item[Pag. 23] CS3014 Estructuras de Datos Avanzadas Signature sort Paso 2 Ya que hay pocos posibles digitos almacenados para un universo muy grande (nlog 1+epsilon n<<2 log2 n), reemplazamos cada bloque por un hash perfecto de O(logn)bits. Notemos que podriamos elegir una constante grande para evitar colisiones con alta probabilidad. Podemos multiplicar cada digito por un enteromy luego obtener las mascaras reducidas de cada uno. Esto tomaraO(1)por digito. Esto preserva la identidad pero no preserva el orden. Paso 3 Notemos que ahora la cantidad de bits usados esO(nlo...
\item[Pag. 24] CS3014 Estructuras de Datos Avanzadas Signature sort Paso 2 Ya que hay pocos posibles digitos almacenados para un universo muy grande (nlog 1+epsilon n<<2 log2 n), reemplazamos cada bloque por un hash perfecto de O(logn)bits. Notemos que podriamos elegir una constante grande para evitar colisiones con alta probabilidad. Podemos multiplicar cada digito por un enteromy luego obtener las mascaras reducidas de cada uno. Esto tomaraO(1)por digito. Esto preserva la identidad pero no preserva el orden. Paso 3 Notemos que ahora la cantidad de bits usados esO(nlo...
\item[Pag. 25] CS3014 Estructuras de Datos Avanzadas Signature sort Paso 2 Ya que hay pocos posibles digitos almacenados para un universo muy grande (nlog 1+epsilon n<<2 log2 n), reemplazamos cada bloque por un hash perfecto de O(logn)bits. Notemos que podriamos elegir una constante grande para evitar colisiones con alta probabilidad. Podemos multiplicar cada digito por un enteromy luego obtener las mascaras reducidas de cada uno. Esto tomaraO(1)por digito. Esto preserva la identidad pero no preserva el orden. Paso 3 Notemos que ahora la cantidad de bits usados esO(nlo...
\item[Pag. 26] CS3014 Estructuras de Datos Avanzadas Signature sort Paso 2 Ya que hay pocos posibles digitos almacenados para un universo muy grande (nlog 1+epsilon n<<2 log2 n), reemplazamos cada bloque por un hash perfecto de O(logn)bits. Notemos que podriamos elegir una constante grande para evitar colisiones con alta probabilidad. Podemos multiplicar cada digito por un enteromy luego obtener las mascaras reducidas de cada uno. Esto tomaraO(1)por digito. Esto preserva la identidad pero no preserva el orden. Paso 3 Notemos que ahora la cantidad de bits usados esO(nlo...
\item[Pag. 27] CS3014 Estructuras de Datos Avanzadas Signature sort Paso 2 Ya que hay pocos posibles digitos almacenados para un universo muy grande (nlog 1+epsilon n<<2 log2 n), reemplazamos cada bloque por un hash perfecto de O(logn)bits. Notemos que podriamos elegir una constante grande para evitar colisiones con alta probabilidad. Podemos multiplicar cada digito por un enteromy luego obtener las mascaras reducidas de cada uno. Esto tomaraO(1)por digito. Esto preserva la identidad pero no preserva el orden. Paso 3 Notemos que ahora la cantidad de bits usados esO(nlo...
\item[Pag. 28] CS3014 Estructuras de Datos Avanzadas Signature sort Paso 2 Ya que hay pocos posibles digitos almacenados para un universo muy grande (nlog 1+epsilon n<<2 log2 n), reemplazamos cada bloque por un hash perfecto de O(logn)bits. Notemos que podriamos elegir una constante grande para evitar colisiones con alta probabilidad. Podemos multiplicar cada digito por un enteromy luego obtener las mascaras reducidas de cada uno. Esto tomaraO(1)por digito. Esto preserva la identidad pero no preserva el orden. Paso 3 Notemos que ahora la cantidad de bits usados esO(nlo...
\item[Pag. 29] CS3014 Estructuras de Datos Avanzadas Signature sort Packed sort Es un algoritmo que nos permite ordenarnenteros debbits enO(n)si el tamano de la palabra cumplew= Omega(blognlog logn). El paso 2 nos dejo con enteros deO(log1+epsilon n)bits, asi que cumplimos la con- dicion sin problemas sobre los hashes/signatures. Por lo tanto, ordenaremos los hashes enO(n). Paso 4 Construimos un trie comprimido sobre los hashes, de manera que el inorder traversal del mismo nos da los hashes ordenados. Esto tomaO(n)de espacio y es posible construirlo enO(n)tambien. Vict...
\item[Pag. 30] CS3014 Estructuras de Datos Avanzadas Signature sort Packed sort Es un algoritmo que nos permite ordenarnenteros debbits enO(n)si el tamano de la palabra cumplew= Omega(blognlog logn). El paso 2 nos dejo con enteros deO(log1+epsilon n)bits, asi que cumplimos la con- dicion sin problemas sobre los hashes/signatures. Por lo tanto, ordenaremos los hashes enO(n). Paso 4 Construimos un trie comprimido sobre los hashes, de manera que el inorder traversal del mismo nos da los hashes ordenados. Esto tomaO(n)de espacio y es posible construirlo enO(n)tambien. Vict...
\item[Pag. 31] CS3014 Estructuras de Datos Avanzadas Signature sort Packed sort Es un algoritmo que nos permite ordenarnenteros debbits enO(n)si el tamano de la palabra cumplew= Omega(blognlog logn). El paso 2 nos dejo con enteros deO(log1+epsilon n)bits, asi que cumplimos la con- dicion sin problemas sobre los hashes/signatures. Por lo tanto, ordenaremos los hashes enO(n). Paso 4 Construimos un trie comprimido sobre los hashes, de manera que el inorder traversal del mismo nos da los hashes ordenados. Esto tomaO(n)de espacio y es posible construirlo enO(n)tambien. Vict...
\item[Pag. 32] CS3014 Estructuras de Datos Avanzadas Signature sort Packed sort Es un algoritmo que nos permite ordenarnenteros debbits enO(n)si el tamano de la palabra cumplew= Omega(blognlog logn). El paso 2 nos dejo con enteros deO(log1+epsilon n)bits, asi que cumplimos la con- dicion sin problemas sobre los hashes/signatures. Por lo tanto, ordenaremos los hashes enO(n). Paso 4 Construimos un trie comprimido sobre los hashes, de manera que el inorder traversal del mismo nos da los hashes ordenados. Esto tomaO(n)de espacio y es posible construirlo enO(n)tambien. Vict...
\item[Pag. 33] CS3014 Estructuras de Datos Avanzadas Signature sort Packed sort Es un algoritmo que nos permite ordenarnenteros debbits enO(n)si el tamano de la palabra cumplew= Omega(blognlog logn). El paso 2 nos dejo con enteros deO(log1+epsilon n)bits, asi que cumplimos la con- dicion sin problemas sobre los hashes/signatures. Por lo tanto, ordenaremos los hashes enO(n). Paso 4 Construimos un trie comprimido sobre los hashes, de manera que el inorder traversal del mismo nos da los hashes ordenados. Esto tomaO(n)de espacio y es posible construirlo enO(n)tambien. Vict...
\item[Pag. 34] CS3014 Estructuras de Datos Avanzadas Signature sort Packed sort Es un algoritmo que nos permite ordenarnenteros debbits enO(n)si el tamano de la palabra cumplew= Omega(blognlog logn). El paso 2 nos dejo con enteros deO(log1+epsilon n)bits, asi que cumplimos la con- dicion sin problemas sobre los hashes/signatures. Por lo tanto, ordenaremos los hashes enO(n). Paso 4 Construimos un trie comprimido sobre los hashes, de manera que el inorder traversal del mismo nos da los hashes ordenados. Esto tomaO(n)de espacio y es posible construirlo enO(n)tambien. Vict...
\item[Pag. 35] CS3014 Estructuras de Datos Avanzadas Signature sort Trie comprimido Tomaremos cada nodo del trie con alfabeto\{0,1, . . . ,B-1\}que no tiene al menos dos hijos almacenados y lo "uniremos" con su padre. Debido a lo anterior, tendremosO(n)nodos en la version comprimida, com- parado con losO(nlog epsilon n)que tendriamos en la version completa. Problematica Tenemos todos los enteros originales ordenados segun sus hashes, asi que en realidad no tenemos ningun orden correcto todavia. Victor Racso Galvan Oyola Universidad de Ingenieria y Tecnologia 11 / 15
\item[Pag. 36] CS3014 Estructuras de Datos Avanzadas Signature sort Trie comprimido Tomaremos cada nodo del trie con alfabeto\{0,1, . . . ,B-1\}que no tiene al menos dos hijos almacenados y lo "uniremos" con su padre. Debido a lo anterior, tendremosO(n)nodos en la version comprimida, com- parado con losO(nlog epsilon n)que tendriamos en la version completa. Problematica Tenemos todos los enteros originales ordenados segun sus hashes, asi que en realidad no tenemos ningun orden correcto todavia. Victor Racso Galvan Oyola Universidad de Ingenieria y Tecnologia 11 / 15
\item[Pag. 37] CS3014 Estructuras de Datos Avanzadas Signature sort Trie comprimido Tomaremos cada nodo del trie con alfabeto\{0,1, . . . ,B-1\}que no tiene al menos dos hijos almacenados y lo "uniremos" con su padre. Debido a lo anterior, tendremosO(n)nodos en la version comprimida, com- parado con losO(nlog epsilon n)que tendriamos en la version completa. Problematica Tenemos todos los enteros originales ordenados segun sus hashes, asi que en realidad no tenemos ningun orden correcto todavia. Victor Racso Galvan Oyola Universidad de Ingenieria y Tecnologia 11 / 15
\item[Pag. 38] CS3014 Estructuras de Datos Avanzadas Signature sort Trie comprimido Tomaremos cada nodo del trie con alfabeto\{0,1, . . . ,B-1\}que no tiene al menos dos hijos almacenados y lo "uniremos" con su padre. Debido a lo anterior, tendremosO(n)nodos en la version comprimida, com- parado con losO(nlog epsilon n)que tendriamos en la version completa. Problematica Tenemos todos los enteros originales ordenados segun sus hashes, asi que en realidad no tenemos ningun orden correcto todavia. Victor Racso Galvan Oyola Universidad de Ingenieria y Tecnologia 11 / 15
\item[Pag. 39] CS3014 Estructuras de Datos Avanzadas Signature sort Trie comprimido Tomaremos cada nodo del trie con alfabeto\{0,1, . . . ,B-1\}que no tiene al menos dos hijos almacenados y lo "uniremos" con su padre. Debido a lo anterior, tendremosO(n)nodos en la version comprimida, com- parado con losO(nlog epsilon n)que tendriamos en la version completa. Problematica Tenemos todos los enteros originales ordenados segun sus hashes, asi que en realidad no tenemos ningun orden correcto todavia. Victor Racso Galvan Oyola Universidad de Ingenieria y Tecnologia 11 / 15
\item[Pag. 40] CS3014 Estructuras de Datos Avanzadas Signature sort Paso 5 Ordenaremos las aristas del trie recursivamente usando (Indice del nodo,Primer digito,Indice del hijo). De esta manera, si el nivel actual tiene enteros dew' bits, estamos reducien- dolos a una tupla de(O(logn),O( w ' logepsilon n ),O(logn)), asi que la cantidad de bits se reduce a w ' logepsilon n +O(logn) Luego de 1 epsilon +1 niveles de recursion se tendra que la cantidad de bits de cada entero esb=O( 1 epsilon logn+ w ' log1+epsilon n ) =O( w lognlog logn ). Sobre este escenario podemos usar...
\item[Pag. 41] CS3014 Estructuras de Datos Avanzadas Signature sort Paso 5 Ordenaremos las aristas del trie recursivamente usando (Indice del nodo,Primer digito,Indice del hijo). De esta manera, si el nivel actual tiene enteros dew' bits, estamos reducien- dolos a una tupla de(O(logn),O( w ' logepsilon n ),O(logn)), asi que la cantidad de bits se reduce a w ' logepsilon n +O(logn) Luego de 1 epsilon +1 niveles de recursion se tendra que la cantidad de bits de cada entero esb=O( 1 epsilon logn+ w ' log1+epsilon n ) =O( w lognlog logn ). Sobre este escenario podemos usar...
\item[Pag. 42] CS3014 Estructuras de Datos Avanzadas Signature sort Paso 5 Ordenaremos las aristas del trie recursivamente usando (Indice del nodo,Primer digito,Indice del hijo). De esta manera, si el nivel actual tiene enteros dew' bits, estamos reducien- dolos a una tupla de(O(logn),O( w ' logepsilon n ),O(logn)), asi que la cantidad de bits se reduce a w ' logepsilon n +O(logn) Luego de 1 epsilon +1 niveles de recursion se tendra que la cantidad de bits de cada entero esb=O( 1 epsilon logn+ w ' log1+epsilon n ) =O( w lognlog logn ). Sobre este escenario podemos usar...
\item[Pag. 43] CS3014 Estructuras de Datos Avanzadas Signature sort Paso 5 Ordenaremos las aristas del trie recursivamente usando (Indice del nodo,Primer digito,Indice del hijo). De esta manera, si el nivel actual tiene enteros dew' bits, estamos reducien- dolos a una tupla de(O(logn),O( w ' logepsilon n ),O(logn)), asi que la cantidad de bits se reduce a w ' logepsilon n +O(logn) Luego de 1 epsilon +1 niveles de recursion se tendra que la cantidad de bits de cada entero esb=O( 1 epsilon logn+ w ' log1+epsilon n ) =O( w lognlog logn ). Sobre este escenario podemos usar...
\item[Pag. 44] CS3014 Estructuras de Datos Avanzadas Signature sort Paso 5 Ordenaremos las aristas del trie recursivamente usando (Indice del nodo,Primer digito,Indice del hijo). De esta manera, si el nivel actual tiene enteros dew' bits, estamos reducien- dolos a una tupla de(O(logn),O( w ' logepsilon n ),O(logn)), asi que la cantidad de bits se reduce a w ' logepsilon n +O(logn) Luego de 1 epsilon +1 niveles de recursion se tendra que la cantidad de bits de cada entero esb=O( 1 epsilon logn+ w ' log1+epsilon n ) =O( w lognlog logn ). Sobre este escenario podemos usar...
\item[Pag. 45] CS3014 Estructuras de Datos Avanzadas Signature sort Paso 5 Ordenaremos las aristas del trie recursivamente usando (Indice del nodo,Primer digito,Indice del hijo). De esta manera, si el nivel actual tiene enteros dew' bits, estamos reducien- dolos a una tupla de(O(logn),O( w ' logepsilon n ),O(logn)), asi que la cantidad de bits se reduce a w ' logepsilon n +O(logn) Luego de 1 epsilon +1 niveles de recursion se tendra que la cantidad de bits de cada entero esb=O( 1 epsilon logn+ w ' log1+epsilon n ) =O( w lognlog logn ). Sobre este escenario podemos usar...
\item[Pag. 46] CS3014 Estructuras de Datos Avanzadas Signature sort Paso 5 Ordenaremos las aristas del trie recursivamente usando (Indice del nodo,Primer digito,Indice del hijo). De esta manera, si el nivel actual tiene enteros dew' bits, estamos reducien- dolos a una tupla de(O(logn),O( w ' logepsilon n ),O(logn)), asi que la cantidad de bits se reduce a w ' logepsilon n +O(logn) Luego de 1 epsilon +1 niveles de recursion se tendra que la cantidad de bits de cada entero esb=O( 1 epsilon logn+ w ' log1+epsilon n ) =O( w lognlog logn ). Sobre este escenario podemos usar...
\item[Pag. 47] CS3014 Estructuras de Datos Avanzadas Signature sort Paso 7 Ya que ahora tendremos los hijos de cada nodo del trie comprimido correc- tamente ordenados segun el valor original de sus datos, podemos ordenar calculando el inorder traversal del trie. En resumen 1. Dividimos cada entero en digitos para poder tener valores mas pequenos. 2. Aplicamos hashing para poder reducir el espacio usado para representar los digitos. 3. Usamos un trie comprimido para "ordenar" los hashes 4. Ordenamos recursivamente las aristas del trie comprimido para ordenarlos segun el...
\item[Pag. 48] CS3014 Estructuras de Datos Avanzadas Signature sort Paso 7 Ya que ahora tendremos los hijos de cada nodo del trie comprimido correc- tamente ordenados segun el valor original de sus datos, podemos ordenar calculando el inorder traversal del trie. En resumen 1. Dividimos cada entero en digitos para poder tener valores mas pequenos. 2. Aplicamos hashing para poder reducir el espacio usado para representar los digitos. 3. Usamos un trie comprimido para "ordenar" los hashes 4. Ordenamos recursivamente las aristas del trie comprimido para ordenarlos segun el...
\item[Pag. 49] CS3014 Estructuras de Datos Avanzadas Signature sort Paso 7 Ya que ahora tendremos los hijos de cada nodo del trie comprimido correc- tamente ordenados segun el valor original de sus datos, podemos ordenar calculando el inorder traversal del trie. En resumen 1. Dividimos cada entero en digitos para poder tener valores mas pequenos. 2. Aplicamos hashing para poder reducir el espacio usado para representar los digitos. 3. Usamos un trie comprimido para "ordenar" los hashes 4. Ordenamos recursivamente las aristas del trie comprimido para ordenarlos segun el...
\item[Pag. 50] CS3014 Estructuras de Datos Avanzadas Signature sort Paso 7 Ya que ahora tendremos los hijos de cada nodo del trie comprimido correc- tamente ordenados segun el valor original de sus datos, podemos ordenar calculando el inorder traversal del trie. En resumen 1. Dividimos cada entero en digitos para poder tener valores mas pequenos. 2. Aplicamos hashing para poder reducir el espacio usado para representar los digitos. 3. Usamos un trie comprimido para "ordenar" los hashes 4. Ordenamos recursivamente las aristas del trie comprimido para ordenarlos segun el...
\item[Pag. 51] CS3014 Estructuras de Datos Avanzadas Signature sort Paso 7 Ya que ahora tendremos los hijos de cada nodo del trie comprimido correc- tamente ordenados segun el valor original de sus datos, podemos ordenar calculando el inorder traversal del trie. En resumen 1. Dividimos cada entero en digitos para poder tener valores mas pequenos. 2. Aplicamos hashing para poder reducir el espacio usado para representar los digitos. 3. Usamos un trie comprimido para "ordenar" los hashes 4. Ordenamos recursivamente las aristas del trie comprimido para ordenarlos segun el...
\item[Pag. 52] CS3014 Estructuras de Datos Avanzadas Signature sort Paso 7 Ya que ahora tendremos los hijos de cada nodo del trie comprimido correc- tamente ordenados segun el valor original de sus datos, podemos ordenar calculando el inorder traversal del trie. En resumen 1. Dividimos cada entero en digitos para poder tener valores mas pequenos. 2. Aplicamos hashing para poder reducir el espacio usado para representar los digitos. 3. Usamos un trie comprimido para "ordenar" los hashes 4. Ordenamos recursivamente las aristas del trie comprimido para ordenarlos segun el...
\item[Pag. 53] CS3014 Estructuras de Datos Avanzadas Signature sort Paso 7 Ya que ahora tendremos los hijos de cada nodo del trie comprimido correc- tamente ordenados segun el valor original de sus datos, podemos ordenar calculando el inorder traversal del trie. En resumen 1. Dividimos cada entero en digitos para poder tener valores mas pequenos. 2. Aplicamos hashing para poder reducir el espacio usado para representar los digitos. 3. Usamos un trie comprimido para "ordenar" los hashes 4. Ordenamos recursivamente las aristas del trie comprimido para ordenarlos segun el...
\item[Pag. 54] CS3014 Estructuras de Datos Avanzadas Signature sort Paso 7 Ya que ahora tendremos los hijos de cada nodo del trie comprimido correc- tamente ordenados segun el valor original de sus datos, podemos ordenar calculando el inorder traversal del trie. En resumen 1. Dividimos cada entero en digitos para poder tener valores mas pequenos. 2. Aplicamos hashing para poder reducir el espacio usado para representar los digitos. 3. Usamos un trie comprimido para "ordenar" los hashes 4. Ordenamos recursivamente las aristas del trie comprimido para ordenarlos segun el...
\item[Pag. 55] CS3014 Estructuras de Datos Avanzadas Resumen de la sesion -Aprendimos la equivalencia entre ordenamiento y colas de prioridad. -Aprendimos algunos resultados de ordenamientos de enteros. -Aprendimos signature sort. Victor Racso Galvan Oyola Universidad de Ingenieria y Tecnologia 14 / 15
\item[Pag. 56] CS3014 Estructuras de Datos Avanzadas Resumen de la sesion -Aprendimos la equivalencia entre ordenamiento y colas de prioridad. -Aprendimos algunos resultados de ordenamientos de enteros. -Aprendimos signature sort. Victor Racso Galvan Oyola Universidad de Ingenieria y Tecnologia 14 / 15
\item[Pag. 57] CS3014 Estructuras de Datos Avanzadas Resumen de la sesion -Aprendimos la equivalencia entre ordenamiento y colas de prioridad. -Aprendimos algunos resultados de ordenamientos de enteros. -Aprendimos signature sort. Victor Racso Galvan Oyola Universidad de Ingenieria y Tecnologia 14 / 15
\item[Pag. 58] CS3014 Estructuras de Datos Avanzadas Resumen de la sesion -Aprendimos la equivalencia entre ordenamiento y colas de prioridad. -Aprendimos algunos resultados de ordenamientos de enteros. -Aprendimos signature sort. Victor Racso Galvan Oyola Universidad de Ingenieria y Tecnologia 14 / 15
\item[Pag. 59] CS3014 Estructuras de Datos Avanzadas Gracias Victor Racso Galvan Oyola Universidad de Ingenieria y Tecnologia 15 / 15
\end{description}
\subsection*{eda\_slides\_06\_integer.pdf (48 paginas)}
\begin{description}[leftmargin=!,labelwidth=4.2em,style=nextline]
\item[Pag. 1] Integer I CS3014 - Estructura de Datos Avanzados Luciano A. Romero Calla lromeroc@utec.edu.pe 2026-1
\item[Pag. 2] Overview 1. Goals 2. Models 3. Predecesor problem 4. van Emde Boas 5. Tree view 6. Saving space 2/19
\item[Pag. 3] Goals 1. Goals 3/19
\item[Pag. 4] Goals Goals - Explore integer data structures to solve the predecessor problem. - Describe two main data structuresvan Emde BoasandY-fasttrees. 4/19
\item[Pag. 5] Goals Goals - Explore integer data structures to solve the predecessor problem. - Describe two main data structuresvan Emde BoasandY-fasttrees. 4/19
\item[Pag. 6] Models 2. Models 5/19
\item[Pag. 7] Models Models word w-bit integer in U=\{0,1,2,3, ...,u=2 w -1\} transdichotomous RAM bridging two worlds machine/problem,w>=lgS word RAM transdichotomous RAM with C-style operations cell probe count \# memory reads \& writes, computation-free, useful for lower bounds cell probe strong transdichotomous RAM word RAM pointer machine BST weak 6/19
\item[Pag. 8] Models Models word w-bit integer in U=\{0,1,2,3, ...,u=2 w -1\} transdichotomous RAM bridging two worlds machine/problem,w>=lgS word RAM transdichotomous RAM with C-style operations cell probe count \# memory reads \& writes, computation-free, useful for lower bounds cell probe strong transdichotomous RAM word RAM pointer machine BST weak 6/19
\item[Pag. 9] Models Models word w-bit integer in U=\{0,1,2,3, ...,u=2 w -1\} transdichotomous RAM bridging two worlds machine/problem,w>=lgS word RAM transdichotomous RAM with C-style operations cell probe count \# memory reads \& writes, computation-free, useful for lower bounds cell probe strong transdichotomous RAM word RAM pointer machine BST weak 6/19
\item[Pag. 10] Models Models word w-bit integer in U=\{0,1,2,3, ...,u=2 w -1\} transdichotomous RAM bridging two worlds machine/problem,w>=lgS word RAM transdichotomous RAM with C-style operations cell probe count \# memory reads \& writes, computation-free, useful for lower bounds cell probe strong transdichotomous RAM word RAM pointer machine BST weak 6/19
\item[Pag. 11] Models Models word w-bit integer in U=\{0,1,2,3, ...,u=2 w -1\} transdichotomous RAM bridging two worlds machine/problem,w>=lgS word RAM transdichotomous RAM with C-style operations cell probe count \# memory reads \& writes, computation-free, useful for lower bounds cell probe strong transdichotomous RAM word RAM pointer machine BST weak 6/19
\item[Pag. 12] Predecesor problem 3. Predecesor problem 7/19
\item[Pag. 13] Predecesor problem Predecesor problem problem maintain a setSofnwords - insert(x in U) - delete(x in S) - predecessor(x in U):max\{y in S|y<x\} - successor(x in U):max\{y in S|y>x\} BST comparison model:O(lgn)/op optimal word RAM - van Emde Boas:O(lgw)/op,Theta(u)space - Y-fast trees:O(lgw)/op,Theta(n)space - fusion trees:O(log w n)/op,Theta(n)space cell probe lower bound O(n polylg n)space,Omega min logw n, lgw lg lgw lg lgn pointer machine - word specified by node pointer - van Emde Boas:O(lg lgu)/op,Theta(u)space - lower bound:Omega(lg lgu)/op,Omega(u)spa...
\item[Pag. 14] Predecesor problem Predecesor problem problem maintain a setSofnwords - insert(x in U) - delete(x in S) - predecessor(x in U):max\{y in S|y<x\} - successor(x in U):max\{y in S|y>x\} BST comparison model:O(lgn)/op optimal word RAM - van Emde Boas:O(lgw)/op,Theta(u)space - Y-fast trees:O(lgw)/op,Theta(n)space - fusion trees:O(log w n)/op,Theta(n)space cell probe lower bound O(n polylg n)space,Omega min logw n, lgw lg lgw lg lgn pointer machine - word specified by node pointer - van Emde Boas:O(lg lgu)/op,Theta(u)space - lower bound:Omega(lg lgu)/op,Omega(u)spa...
\item[Pag. 15] Predecesor problem Predecesor problem problem maintain a setSofnwords - insert(x in U) - delete(x in S) - predecessor(x in U):max\{y in S|y<x\} - successor(x in U):max\{y in S|y>x\} BST comparison model:O(lgn)/op optimal word RAM - van Emde Boas:O(lgw)/op,Theta(u)space - Y-fast trees:O(lgw)/op,Theta(n)space - fusion trees:O(log w n)/op,Theta(n)space cell probe lower bound O(n polylg n)space,Omega min logw n, lgw lg lgw lg lgn pointer machine - word specified by node pointer - van Emde Boas:O(lg lgu)/op,Theta(u)space - lower bound:Omega(lg lgu)/op,Omega(u)spa...
\item[Pag. 16] Predecesor problem Predecesor problem problem maintain a setSofnwords - insert(x in U) - delete(x in S) - predecessor(x in U):max\{y in S|y<x\} - successor(x in U):max\{y in S|y>x\} BST comparison model:O(lgn)/op optimal word RAM - van Emde Boas:O(lgw)/op,Theta(u)space - Y-fast trees:O(lgw)/op,Theta(n)space - fusion trees:O(log w n)/op,Theta(n)space cell probe lower bound O(n polylg n)space,Omega min logw n, lgw lg lgw lg lgn pointer machine - word specified by node pointer - van Emde Boas:O(lg lgu)/op,Theta(u)space - lower bound:Omega(lg lgu)/op,Omega(u)spa...
\item[Pag. 17] Predecesor problem Predecesor problem problem maintain a setSofnwords - insert(x in U) - delete(x in S) - predecessor(x in U):max\{y in S|y<x\} - successor(x in U):max\{y in S|y>x\} BST comparison model:O(lgn)/op optimal word RAM - van Emde Boas:O(lgw)/op,Theta(u)space - Y-fast trees:O(lgw)/op,Theta(n)space - fusion trees:O(log w n)/op,Theta(n)space cell probe lower bound O(n polylg n)space,Omega min logw n, lgw lg lgw lg lgn pointer machine - word specified by node pointer - van Emde Boas:O(lg lgu)/op,Theta(u)space - lower bound:Omega(lg lgu)/op,Omega(u)spa...
\item[Pag. 18] van Emde Boas 4. van Emde Boas 4.1 Definition 4.2 Operations 9/19
\item[Pag. 19] van Emde Boas Definition Definition idea T(u) =T( sqrtu) +O(1) hierarchical coordinates wordx= c,i recursive vEBVof sizeu - V.cluster[i] =vEB of size sqrtu - V.summary[i] =vEB of size sqrtu - V.min=minimum element inV - V.max=maximum element inV 10/19
\item[Pag. 20] van Emde Boas Definition Definition idea T(u) =T( sqrtu) +O(1) hierarchical coordinates wordx= c,i recursive vEBVof sizeu - V.cluster[i] =vEB of size sqrtu - V.summary[i] =vEB of size sqrtu - V.min=minimum element inV - V.max=maximum element inV 10/19
\item[Pag. 21] van Emde Boas Definition Definition idea T(u) =T( sqrtu) +O(1) hierarchical coordinates wordx= c,i recursive vEBVof sizeu - V.cluster[i] =vEB of size sqrtu - V.summary[i] =vEB of size sqrtu - V.min=minimum element inV - V.max=maximum element inV 10/19
\item[Pag. 22] van Emde Boas Definition Definition idea T(u) =T( sqrtu) +O(1) hierarchical coordinates wordx= c,i recursive vEBVof sizeu - V.cluster[i] =vEB of size sqrtu - V.summary[i] =vEB of size sqrtu - V.min=minimum element inV - V.max=maximum element inV 10/19
\item[Pag. 23] van Emde Boas Definition Definition idea T(u) =T( sqrtu) +O(1) hierarchical coordinates wordx= c,i recursive vEBVof sizeu - V.cluster[i] =vEB of size sqrtu - V.summary[i] =vEB of size sqrtu - V.min=minimum element inV - V.max=maximum element inV 10/19
\item[Pag. 24] van Emde Boas Definition Definition idea T(u) =T( sqrtu) +O(1) hierarchical coordinates wordx= c,i recursive vEBVof sizeu - V.cluster[i] =vEB of size sqrtu - V.summary[i] =vEB of size sqrtu - V.min=minimum element inV - V.max=maximum element inV 10/19
\item[Pag. 25] van Emde Boas Definition Definition Figure 1:van Emde Boas structure 1 1Demaine 2021. 11/19
\item[Pag. 26] van Emde Boas Operations Operations Successor(V,x= c,i ) - ifx<V.min returnV.min - ifi<V.cluster[c].max return c,Successor(V.cluster[c],i) - elsec ' =Successor(V.summary,c) return c ',V.cluster[c '].min 12/19
\item[Pag. 27] Tree view 5. Tree view 5.1 Indirection 13/19
\item[Pag. 28] Tree view Tree view Figure 2:van Emde Boas tree view 2 2Demaine 2021. 14/19
\item[Pag. 29] Tree view Indirection Indirection - fromO(lgw)query,O(w)update DS - reduce toO(lgw)update: splitnelements into chuncks of sizeTheta(w) 15/19
\item[Pag. 30] Tree view Indirection Indirection - fromO(lgw)query,O(w)update DS - reduce toO(lgw)update: splitnelements into chuncks of sizeTheta(w) 15/19
\item[Pag. 31] Saving space 6. Saving space 6.1 X-fast trees 6.2 Y-fast trees 16/19
\item[Pag. 32] Saving space Saving space - don't store empty clusters vEB - V.clusters = hash table - charge each table entry to min in child - space boundTheta(n)via indirection 17/19
\item[Pag. 33] Saving space Saving space - don't store empty clusters vEB - V.clusters = hash table - charge each table entry to min in child - space boundTheta(n)via indirection 17/19
\item[Pag. 34] Saving space Saving space - don't store empty clusters vEB - V.clusters = hash table - charge each table entry to min in child - space boundTheta(n)via indirection 17/19
\item[Pag. 35] Saving space Saving space - don't store empty clusters vEB - V.clusters = hash table - charge each table entry to min in child - space boundTheta(n)via indirection 17/19
\item[Pag. 36] Saving space X-fast trees X-fast trees - don't store 0s - store a hash table of root-to-1 path - O(lgw)query binary search - Theta(w)update - Theta(nw)space 18/19
\item[Pag. 37] Saving space X-fast trees X-fast trees - don't store 0s - store a hash table of root-to-1 path - O(lgw)query binary search - Theta(w)update - Theta(nw)space 18/19
\item[Pag. 38] Saving space X-fast trees X-fast trees - don't store 0s - store a hash table of root-to-1 path - O(lgw)query binary search - Theta(w)update - Theta(nw)space 18/19
\item[Pag. 39] Saving space X-fast trees X-fast trees - don't store 0s - store a hash table of root-to-1 path - O(lgw)query binary search - Theta(w)update - Theta(nw)space 18/19
\item[Pag. 40] Saving space X-fast trees X-fast trees - don't store 0s - store a hash table of root-to-1 path - O(lgw)query binary search - Theta(w)update - Theta(nw)space 18/19
\item[Pag. 41] Saving space Y-fast trees Y-fast trees - = X-fast trees - + indirection - O(lgw)query binary search - O(lgw)update amortized - O(lgw)space 19/19
\item[Pag. 42] Saving space Y-fast trees Y-fast trees - = X-fast trees - + indirection - O(lgw)query binary search - O(lgw)update amortized - O(lgw)space 19/19
\item[Pag. 43] Saving space Y-fast trees Y-fast trees - = X-fast trees - + indirection - O(lgw)query binary search - O(lgw)update amortized - O(lgw)space 19/19
\item[Pag. 44] Saving space Y-fast trees Y-fast trees - = X-fast trees - + indirection - O(lgw)query binary search - O(lgw)update amortized - O(lgw)space 19/19
\item[Pag. 45] Saving space Y-fast trees Y-fast trees - = X-fast trees - + indirection - O(lgw)query binary search - O(lgw)update amortized - O(lgw)space 19/19
\item[Pag. 46] Thanks
\item[Pag. 47] References References Demaine, Erik (2021).6.851: Advanced Data Structures (Spring'21).url: https://courses. csail.mit.edu/6.851/spring21/. 21/19
\item[Pag. 48] Acknowledgements Acknowledgements The course slides are based on the lectures from previous editions and on similar lectures elsewhere. List of credits: Erick Demaine, Keith Schwarz 22/19
\end{description}
